\section{Problem, scope and conventions}\label{sec:scope}

The design question is the following: given a smooth target interface inside a
cylinder, which admissible acoustic source distributions can form it and
maintain it with specified accuracy and stability? The target is initially an
arbitrary surface. Its compatibility with the fluid volume and boundaries,
the required stress, the range of realizable acoustic loads, and its dynamic
stability are separate conditions. Cartesian optical surfaces enter only after
these general conditions have been formulated.

Existing observations establish deformation by radiation pressure and the
importance of the interaction between the wave and its evolving interface
\cite{bertin2012,chesneau2022}. Time-dependent deformations have also been
measured under localized ultrasound \cite{sisombat2023}. These results motivate
the present formulation, but do not establish arbitrary surface controllability
or optical precision for an unspecified apparatus.

\subsection{Cylindrical geometry and data left as arguments}
Let $t$ denote physical time, and let $T_f>0$ denote a chosen formation horizon.
The position vector is $\bm x=(x_1,x_2,z)\in\R^3$. Define a reference cylinder
of radius $R_c>0$ and height $H_c=z_t-z_b>0$ by
\begin{equation}
 \mathcal D_c=\{\bm x:x_1^2+x_2^2<R_c^2,\ z_b<z<z_t\}.
 \label{eq:cylinder}
\end{equation}
Here $z_b$ and $z_t$ are the bottom and top coordinates. The cylinder fixes
the container geometry, not a symmetry of the state. An interface, source
field or perturbation may depend on all three coordinates. A connected fluid region
$\Omega^-(t)$ is surrounded, at least across its optical interface, by a second
fluid region $\Omega^+(t)$. The superscript $-$ labels the lens material and $+$
the adjoining fluid; either may be the incident optical medium. A gas is
included as a possible adjoining fluid, without replacing it by a vacuum or a
pressure-release boundary at this stage.

Let $\Gamma(t)$ be the deformable fluid--fluid interface and
$\Gamma_{\rm opt}(t)\subseteq\Gamma(t)$ its clear optical part. Other boundaries
may support the liquid, radiate acoustic energy, or couple to elastic solids.
The actuator boundary, when actuation is imposed on a boundary, is denoted by
$\Gamma_a$. The contact line, when one exists, is
$\mathcal C(t)=\partial\Gamma(t)$ on a supporting solid. The two phases
partition the fluid-filled part of the cylinder; additional immersed bodies
require their geometry and boundary conditions as explicit data. Closed free
interfaces have no contact line. A graph over a circular disk is one admissible
representation, not a requirement; overhangs need a general surface description.

When acoustic motion is prescribed on a wall, $\partial\mathcal D_c$ is its
reference or cycle-mean position. The full waveform problem uses the moving
wall, and its linearization transfers the condition to that reference surface.
An exactly fixed impermeable wall cannot also have a nonzero normal velocity.
Any open boundary or volume source must declare its mass and energy exchange.

The unit normal $\nn$ points from $\Omega^-$ to $\Omega^+$. Define the identity
tensor $\Id$, tangential projector $\Pg=\Id-\nn\otimes\nn$, surface gradient
$\Gs=\Pg\nabla$, and signed total curvature
$\kappa=\Gs\cdot\nn$. Thus a spherical lens-fluid region has positive
curvature. The symbol $\otimes$ denotes a tensor product. For a quantity with
two interface traces, define its jump as
$\jump{f}=f^+-f^-$. These choices fix the signs of all interfacial forces.

The supplied physical data must include the fluid constitutive laws,
body force, external boundaries, wetting law or pinning
constraint, conserved fluid masses, actuator response and admissible command
set. A device-independent formulation leaves those data as arguments of the
problem; it cannot eliminate their influence on the answer. In particular,
optical conjugates alone cannot determine electrical voltages. No particular
aperture, fill height, immersion fluid, source count or operating frequency is
assigned. Arbitrary sources means arbitrary functions within a declared
physical admissible set, not independent prescription of pressure and velocity
at every interface point.

\subsection{Three meanings of a stable optical result}
Let $\Gamma_\star$ denote a desired interface. We distinguish:
\begin{enumerate}
\item \emph{Transient formation}: the trajectory approaches $\Gamma_\star$
within a finite time, possibly with nonzero fluid velocity or acoustic energy.
\item \emph{Driven operation}: constant, periodic, or feedback-controlled
actuation maintains acceptable optics and rejects specified disturbances.
\item \emph{Unforced retention}: the shape remains acceptable after the sound,
associated flows, and thermal transients have decayed.
\end{enumerate}
A pulse means actuation of finite temporal support, with a specified waveform
and bandwidth. It does not imply an ideal Dirac impulse. A periodically repeated
pulse train is continued actuation. Solidification is a change of constitutive
law and belongs to a distinct retention problem.

\begin{figure}[htbp]
\centering
\begin{tikzpicture}[>=Latex,node distance=7mm and 8mm,
 box/.style={draw=blue!40!black,rounded corners,align=center,
 text width=42mm,minimum height=13mm,font=\small}]
 \node[box] (opt) {Admissible target\ and fluid boundaries};
 \node[box,right=of opt] (mech) {Volume, boundaries and\ required fluid traction};
 \node[box,right=of mech] (wave) {Distributed sources and\ acoustic momentum transfer};
 \node[box,below=of mech] (dyn) {Formation trajectory,\ stability and observation};
 \draw[->] (opt)--(mech);
 \draw[<->] (mech)--(wave);
 \draw[->] (wave.south)|-(dyn.east);
 \draw[->] (dyn.west)-|(opt.south);
\end{tikzpicture}
\caption{Dependencies of the inverse problem. Changes of shape alter the acoustic
field; optical performance must be assessed along the resulting physical state.}
\label{fig:dependencies}
\end{figure}

\subsection{Model hierarchy and mathematical status}
The reference continuum description assumes immiscible phases, a regular sharp
interface, and no mass transfer through that interface during shaping. A
Newtonian stress law is supplied explicitly, while non-Newtonian rheology,
solidification, cavitation and topology changes require additional equations.
These are physical modeling choices, not assumptions about a particular
container. Three conditional reductions are identified where used:
\begin{enumerate}[label=(\roman*)]
\item geometric optics in homogeneous isotropic optical media;
\item small acoustic perturbations with a justified averaging timescale;
\item quiescent, constant-density, constant-surface-tension equilibria for an
explicit shape-to-traction formula.
\end{enumerate}
The dynamic formulation retains fluid inertia. Creeping flow appears only as an
explicit limiting case. A stationary interface may coexist with streaming;
such a state requires a steady-flow solve and is not a quiescent equilibrium.

Derivatives below are conditional on a regular local solution and differentiable
state and geometry maps. Smooth interface patches are taken at least $C^3$ for
the displayed curvature variations. We neither prove global well-posedness of
the moving-interface problem nor assert existence of a solution for every
target. Propositions are stated with the assumptions needed for their proofs.
The optimization objectives and design workflow are a proposed synthesis of
these equations. Explicit weak forms, source adjoints and conditional proofs
are given for identified subproblems; the full thermoviscous moving-interface
system is a reference formulation, not a proved globally solvable inverse.
Numerical continuation is neither an assumption nor a step required by the
inverse equations. Numerical experiments will test this formulation.

\subsection{Notation used throughout}
The table defines the principal quantities before their subsequent use. More
specialized symbols are introduced immediately before their equations. All
physical lengths, times and stresses use SI units unless nondimensionalization
is stated.
\begin{longtable}{@{}p{29mm}p{112mm}@{}}
\toprule Symbol & Meaning and units \\\midrule\endhead
$\rho_\ell,\bm v_\ell,p_\ell$ & Instantaneous density, velocity and pressure in
phase $\ell\in\{-,+\}$; $\mathrm{kg\,m^{-3}}$, $\mathrm{m\,s^{-1}}$, $\mathrm{Pa}$.\\
$\bm T_\ell,\bm\tau_\ell$ & Total Cauchy stress and viscous stress; $\mathrm{Pa}$.\\
$\mu_\ell,\zeta_\ell$ & Shear and bulk viscosities; $\mathrm{Pa\,s}$.\\
$\Theta_\ell,e_\ell,k_\ell$ & Temperature, specific internal energy and thermal
conductivity; $\mathrm K$, $\mathrm{J\,kg^{-1}}$, $\mathrm{W\,m^{-1}\,K^{-1}}$.\\
$\sigma,\Phi$ & Surface tension and conservative body-force potential per unit
mass; $\mathrm{N\,m^{-1}}$, $\mathrm{m^2\,s^{-2}}$. The body acceleration is $-\nabla\Phi$.\\
$\lambda_V$ & Pressure multiplier imposing lens volume in an incompressible
reduction; $\mathrm{Pa}$.\\
$\bm U_\ell,p^h_\ell$ & Slow hydrodynamic velocity and pressure in an averaged
description; $\mathrm{m\,s^{-1}}$, $\mathrm{Pa}$.\\
$P_\ell,\bm V_\ell$ & Complex peak acoustic pressure and velocity amplitudes;
$\mathrm{Pa}$ and $\mathrm{m\,s^{-1}}$. They are distinct from the slow fields.\\
$\bm R_\ell,\bm t_{\rm ac},\Pi$ & Acoustic momentum-flux tensor, effective
interfacial acoustic traction, and its outward normal component; $\mathrm{Pa}$.\\
$n_o,n_i,\lambda_{\rm opt}$ & Incident and transmitted optical indices
(dimensionless), and vacuum optical wavelength ($\mathrm m$).\\
$\bm O,\bm I$ & Object and image points; the object subscript is the letter $o$.\\
$\eta,\bm w$ & Normal surface displacement ($\mathrm m$) and complex peak
mechanical emitter velocities ($\mathrm{m\,s^{-1}}$ per component).\\
$\bm a,\bm y$ & General actuator command field and full physical state; each
may be an infinite-dimensional collection of fields. Units depend on the
declared actuator and constitutive model.\\
\bottomrule
\end{longtable}
Bold symbols denote vectors or tensors. The dagger $\dagger$ denotes conjugate
transpose, the superscript $T$ ordinary transpose, and an asterisk complex
conjugation. Inner products are real duality pairings unless explicitly complex;
$\Rea$ and $\Ima$ denote real and imaginary parts.
